\section{Introduction}
\label{sec:introduction}

Cutting a field theory into subregions creates artificial boundaries whose data cannot be reidentified by a formal equality alone. A useful gluing prescription must specify an action principle, balance the outward responses on the two faces, preserve the appropriate covariant phase-space form, and recover both the decoupled and uncut theories in controlled limits. In gauge theory it must also retain independent gauge transformations at the cut until the presymplectic quotient is taken.

We implement this idea with a quadratic interface interaction. At finite coupling the mismatch of the two dressed boundary traces is dynamical and sources equal and opposite outward responses. The weak-coupling endpoint is the direct sum of the region theories, while bounded energy at strong coupling suppresses the mismatch and produces the uncut classical theory. Because all examples are Gaussian, the finite-coupling modes and their CPS normalization can be found explicitly.

Section~\ref{sec:gluing-formalism} gives the general interface variation and CPS construction. Section~\ref{sec:interval-scalar} solves the scalar interval, including its two endpoint theories. Section~\ref{sec:scalar-geometries} treats global $\mathrm{AdS}_2$ and a four-quadrant interface network, thereby separating curvature from nonseparability. Section~\ref{sec:maxwell} introduces the transition field required for gauge-invariant Maxwell gluing and derives both the photon tower and holonomy sector. Section~\ref{sec:numerics} then realizes every oscillator model in the article by mode truncation and matrix diagonalization, including the scalar and Maxwell Dirichlet interpolations of Appendix~\ref{app:dirichlet}. The numerical analysis distinguishes scalar stiffness defects from Maxwell kinetic defects and records the exact claim boundary of the cutoff calculation.
